\documentclass[12pt,a4paper]{article}

% ---- Encoding & Font ----
\usepackage{iftex}
\ifPDFTeX
  \usepackage[utf8]{inputenc}
  \usepackage[T1]{fontenc}
\else
  \usepackage{fontspec}  % XeTeX/LuaTeX (tectonic): Unicode nativo (·, —, ¿, ª, §)
\fi
\usepackage[spanish]{babel}
\usepackage{csquotes}

% ---- Page layout ----
\usepackage[top=2.5cm, bottom=2.5cm, left=2.5cm, right=2.5cm]{geometry}
\usepackage{setspace}
\onehalfspacing

% ---- Math ----
\usepackage{amsmath, amssymb, amsthm}
\usepackage{mathtools}
\usepackage{physics}
\usepackage{esint}  % \oiint

% ---- Graphics ----
\usepackage{graphicx}
\usepackage{tikz}
\usepackage{float}
\usepackage{caption}

% ---- Tables ----
\usepackage{array}
\usepackage{booktabs}
\usepackage{tabularx}
\usepackage{colortbl}

% ---- Colours ----
\usepackage{xcolor}
\definecolor{notebg}{HTML}{FFF3CD}
\definecolor{noteborder}{HTML}{FFC107}
\definecolor{boxred}{HTML}{F8D7DA}
\definecolor{boxredborder}{HTML}{F5C6CB}

% ---- Boxes ----
\usepackage{mdframed}
\usepackage{tcolorbox}
\tcbuselibrary{most}

\newtcolorbox{keybox}[1][]{
  colback=blue!5,
  colframe=blue!70!black,
  arc=4pt,
  boxrule=1pt,
  fonttitle=\bfseries,
  title={#1}
}

\newtcolorbox{warnbox}[1][]{
  colback=boxred,
  colframe=boxredborder,
  coltitle=black!75,
  arc=4pt,
  boxrule=1pt,
  fonttitle=\bfseries,
  title={#1}
}

\newtcolorbox{notebox}[1][]{
  colback=notebg,
  colframe=noteborder,
  coltitle=black!75,
  arc=4pt,
  boxrule=1pt,
  fonttitle=\bfseries,
  title={#1}
}

% ---- Hyperlinks & TOC ----
\usepackage[hidelinks]{hyperref}
\usepackage{bookmark}

% ---- Enumerate ----
\usepackage{enumitem}

% ---- Miscellaneous ----
\usepackage{icomma}
\usepackage{siunitx}
\usepackage{textcomp}
\usepackage[bottom]{footmisc}

% ---- Title ----
\title{\textbf{Clase 13: Cierre del Modelo de Drude y Circuitos \\
        de Corriente Directa}\\[0.3em]
        \large{Conductividad microscópica, leyes de Kirchhoff, baterías
              y FEM, asociación de resistencias y potencia}}
\author{Transcripto y expandido --- Curso Física III (OpenFING)}
\date{}

\begin{document}

\maketitle
\thispagestyle{empty}
\newpage

\tableofcontents
\newpage

% ===================================================================
\section{Cierre del modelo de Drude}
% ===================================================================

Modelo clásico \enquote{de juguete} que justifica la ley de Ohm desde el
movimiento microscópico. La descripción rigurosa es cuántica, pero da la
imagen cualitativa correcta.

\subsection{Repaso de las hipótesis}

\begin{enumerate}[nosep]
  \item Vale la mecánica clásica.
  \item Los electrones no interactúan entre sí (correcta pero difícil de
        justificar).
  \item Movimiento libre con choques cada tiempo medio $\tau$.
  \item Tras el choque, el electrón pierde memoria (arranca en dirección
        aleatoria).
  \item Velocidad promedio $\gg$ velocidad de deriva.
\end{enumerate}

\begin{notebox}[Números (hipótesis 5)]
  Cobre: velocidad promedio $\sim 1{,}6\times10^{6}$~m/s vs.\ deriva
  $\sim 14$~cm/hora. Escalas muy separadas.
\end{notebox}

\subsection{Deducción de la velocidad de deriva}

Electrón que arranca en $t=0$ con velocidad aleatoria $\mathbf{v}_0$;
única fuerza entre choques:
\[
\mathbf{F} = -e\,\mathbf{E} = m\,\mathbf{a}
\;\Longrightarrow\;
\mathbf{a} = -\frac{e}{m}\mathbf{E}
\;\Longrightarrow\;
\mathbf{v}(t) = \mathbf{v}_0 - \frac{e}{m}\mathbf{E}\, t.
\]

Promediando: $\langle\mathbf{v}_0\rangle = 0$ (aleatoria) y
$\langle t\rangle = \tau$:

\begin{keybox}[Velocidad de deriva]
  \[
  \boxed{v_d = \frac{e\,E}{m}\,\tau}
  \]
  El grueso $\mathbf{v}_0$ se cancela; sobrevive el arrastre según
  $\mathbf{E}$. $\tau$ casi no depende de $E$ (propiedad del material).
\end{keybox}

\subsection{Ley de Ohm y conductividad microscópica}

Con $J = e\,n\,v_d$:
\[
J = e\,n\cdot\frac{e\,E}{m}\,\tau = \frac{e^2 n\,\tau}{m}\,E,
\]
que es $J = \sigma E$ con

\begin{keybox}[Conductividad de Drude]
  \[
  \boxed{\sigma = \frac{e^2 n\,\tau}{m}}
  \]
  La proporcionalidad $J\propto E$ vale porque $\tau$ no depende de $E$.
\end{keybox}

\subsection{Con qué chocan realmente los electrones}

Drude imaginaba choques con todos los iones; en realidad los electrones
atraviesan la red casi \textbf{transparente} y solo chocan con
\textbf{impurezas} y \textbf{deformaciones} (efecto cuántico). Por eso
chocan mucho menos de lo esperado.

\begin{notebox}[Temperatura]
  Al bajar $T$ disminuyen deformaciones y defectos, y baja la
  resistividad ($\rho \propto T$ a bajas $T$ para un gas de electrones).
\end{notebox}

\subsection{Efecto Joule}

En promedio los electrones no ganan energía: el trabajo del campo se
entrega en cada choque a las \textbf{vibraciones de la red}
(temperatura). La resistencia se calienta: \textbf{efecto Joule}.

% ===================================================================
\section{Circuitos de corriente directa}
% ===================================================================

\textbf{Corriente directa:} el sentido no cambia (a diferencia de la
alterna, sinusoidal).

\subsection{Aproximación de circuito ideal}

\begin{itemize}[nosep]
  \item \textbf{Cables ideales:} resistencia despreciable frente a los
        resistores.
  \item \textbf{Corriente estacionaria:} tras un transitorio $\sim
        10^{-9}$~s, la corriente en cada punto no depende del tiempo.
  \item \textbf{Corriente igual} a lo largo de un tramo (entre nudos).
  \item \textbf{No se acumula carga:} todas las partes permanecen
        neutras.
\end{itemize}

\begin{warnbox}[Cortocircuito]
  Unir la batería solo con un cable da resistencia diminuta y corriente
  enorme (saltan las llaves térmicas).
\end{warnbox}

\subsection{Primera ley de Kirchhoff (nudos)}

\begin{keybox}[Ley de los nudos]
  \[
  \boxed{I_1 = I_2 + I_3}
  \]
  Toda la corriente que entra a un nudo sale (conservación y no
  acumulación de carga).
\end{keybox}

% ===================================================================
\section{Baterías y fuerza electromotriz}
% ===================================================================

\subsection{Fuentes de energía eléctrica}

Una batería transforma energía en energía eléctrica:
\begin{itemize}[nosep]
  \item \textbf{Pila:} química $\to$ eléctrica.
  \item \textbf{Generador (represa):} mecánica $\to$ eléctrica (efecto
        dínamo).
  \item \textbf{Térmica/nuclear:} combustión/reacción nuclear $\to$
        calor $\to$ turbina $\to$ generador.
  \item \textbf{Solar:} reacción química inducida por luz. \textbf{Eólica:} viento $\to$ turbina.
\end{itemize}
Los generadores realistas dan corriente alterna.

\subsection{Fuerza electromotriz (FEM)}

\begin{keybox}[FEM]
  \[
  \varepsilon = \frac{\text{trabajo entregado al circuito}}{\text{carga que circuló}}
  \]
  Constante, propiedad intrínseca de la batería ($\sim 1{,}5$~V las
  comunes).
\end{keybox}

\begin{warnbox}[Nombre infeliz]
  \enquote{Fuerza electromotriz} \textbf{no es una fuerza} (es energía por
  unidad de carga; nombre histórico), y no coincide exactamente con la FEM
  del curso de Electromagnetismo.
\end{warnbox}

\subsection{Batería ideal y resistencia interna}

\textbf{Ideal:} $|\Delta V| = \varepsilon$. \textbf{Real:} pérdidas por
efecto Joule (resistencia interna $r$):
\[
\Delta V = \varepsilon - r\,I < \varepsilon,
\]
modelada como pila ideal $\varepsilon$ en serie con $r$.

\begin{notebox}[Convención]
  Salvo aviso, todos los dispositivos son ideales.
\end{notebox}

% ===================================================================
\section{Análisis de circuitos: ley de mallas}
% ===================================================================

\subsection{Circuito simple}

Batería $\varepsilon$ con resistencia $R$. El cable ideal ($R=0$) no cae
potencial:
\[
\varepsilon = R\,I \;\Longrightarrow\; \boxed{I = \frac{\varepsilon}{R}}.
\]

\subsection{Segunda ley de Kirchhoff}

Dando la vuelta: subir $\varepsilon$ al atravesar la batería, bajar $R\,I$
al atravesar la resistencia a favor de la corriente:

\begin{keybox}[Ley de mallas]
  \[
  \boxed{\sum_{\text{malla}} \Delta V = 0}
  \]
\end{keybox}

\begin{itemize}[nosep]
  \item Dos resistencias en serie: $I = \varepsilon/(R_1+R_2)$.
  \item Pila con resistencia interna en $R$: $I = \varepsilon/(R+r)$,
        $\Delta V_R = \varepsilon R/(R+r)$; ideal si $R \gg r$.
\end{itemize}

\subsection{Circuito de dos mallas}

Con varias FEM el sentido no es evidente. Estrategia:
\begin{enumerate}[nosep]
  \item Inventar un sentido para cada $I_i$ (si sale negativo, es al
        revés).
  \item Ley de mallas para cada malla independiente (resistencia a favor:
        $-RI$; en contra: $+RI$; FEM según su flecha).
  \item Ley de nudos ($I_3 = I_1 + I_2$).
\end{enumerate}
Queda un sistema lineal $3\times3$. Los signos dependen de las FEM
relativas ($I \propto \varepsilon_1 - \varepsilon_2$).

% ===================================================================
\section{Resistencias en serie y en paralelo}
% ===================================================================

\textbf{Serie} (misma $I$):
\[
V_A - V_B = (R_1+R_2)\,I \;\Longrightarrow\; \boxed{R_{\text{eq}} = R_1 + R_2}.
\]

\textbf{Paralelo} (mismo $\Delta V$, $I = I_1+I_2$):
\[
\boxed{\frac{1}{R_{\text{eq}}} = \frac{1}{R_1} + \frac{1}{R_2}}.
\]

\begin{warnbox}[Al revés que los condensadores]
  En condensadores la suma directa era en paralelo; en resistencias, en
  serie.
\end{warnbox}

% ===================================================================
\section{Transferencia de energía en un circuito}
% ===================================================================

Para un dispositivo entre $A$ y $B$ con corriente $I$ y $\Delta V_{AB}$:
$dW = \Delta V_{AB}\,dQ$, y dividiendo por $dt$:

\begin{keybox}[Potencia entregada]
  \[
  \boxed{P = \Delta V_{AB}\, I}
  \]
\end{keybox}

\begin{center}
\begin{tabular}{lcc}
\toprule
\textbf{Dispositivo} & $\Delta V_{AB}$ & \textbf{Potencia} \\
\midrule
Batería ideal & $\varepsilon$ & $P = \varepsilon I$ (entrega) \\
Resistor & $-R\,I$ & $P = -R\,I^2$ (disipa $R\,I^2$) \\
Condensador & $\pm Q/C$ & almacena o entrega \\
\bottomrule
\end{tabular}
\end{center}

\begin{notebox}[Próxima clase]
  Experimento con circuitos de resistencias y condensadores.
\end{notebox}

\end{document}
